\begin{statementbox}
设双曲线 $x^2/a^2-y^2/b^2=1$，$a,b>0$，$c^2=a^2+b^2$。直线过所选焦点 $F(\pm c,0)$，并有两个不同的有限交点 $A,B$。这里不要求焦点在线段 $AB$ 内。
\begin{center}\begin{tikzpicture}[font=\small,>=stealth]
\draw[rounded corners=3pt,gray!45] (0,0) rectangle (14.2,12.2);
\node at (7.1,11.65) {先看站位：双曲线的“过焦点”，说的是整条直线};\node at (3.6,10.45) {同支：焦点在中间};
\begin{scope}[shift={(3.6,6.1)},scale=0.28]
\begin{scope}\clip (-10,-9) rectangle (10,9);
\draw[teal!75!black,thick,samples=110,domain=-1.9:1.9,variable=\t] plot ({3*cosh(\t)},{4*sinh(\t)});
\draw[teal!75!black,thick,samples=110,domain=-1.9:1.9,variable=\t] plot ({-3*cosh(\t)},{4*sinh(\t)});
\end{scope}
\draw[->,gray] (-10.2,0)--(10.3,0);
\draw[orange!85!black,very thick] (5.0,0)--(4.2000000000,-3.9191835885);\draw[blue!75!black,very thick] (5.0,0)--(6.6000000000,7.8383671769);
\fill (6.6000000000,7.8383671769) circle(6pt) node[above right=7pt] {$A$};\fill (4.2000000000,-3.9191835885) circle(6pt) node[below=8pt] {$B$};
\fill (5.0,0) circle(6pt) node[below right=9pt] {$F$};
\end{scope}
\node at (3.6,2.6) {$A-F-B$：两条反向射线};
\node at (3.6,1.55) {倒数\textbf{相加}};
\node at (10.6,10.45) {异支：焦点在外面};
\begin{scope}[shift={(10.6,6.1)},scale=0.28]
\begin{scope}\clip (-10,-9) rectangle (10,9);
\draw[teal!75!black,thick,samples=110,domain=-1.9:1.9,variable=\t] plot ({3*cosh(\t)},{4*sinh(\t)});
\draw[teal!75!black,thick,samples=110,domain=-1.9:1.9,variable=\t] plot ({-3*cosh(\t)},{4*sinh(\t)});
\end{scope}
\draw[->,gray] (-10.2,0)--(10.3,0);
\draw[orange!85!black,very thick] (5.0,0)--(-5.4173558340,6.0144631950);\draw[blue!75!black,very thick] (5.0,0)--(3.1096635263,1.0913862719);
\fill (3.1096635263,1.0913862719) circle(6pt) node[above right=7pt] {$A$};\fill (-5.4173558340,6.0144631950) circle(6pt) node[left=14pt] {$B$};
\fill (5.0,0) circle(6pt) node[below right=9pt] {$F$};
\end{scope}
\node at (10.6,2.6) {$F-A-B$：同一条射线};
\node at (10.6,1.55) {倒数\textbf{相减}，取绝对值};
\node at (7.1,.55) {两图同为 $x^2/9-y^2/16=1$、同尺度；$F=(5,0)$。};
\end{tikzpicture}
\end{center}
\[
 \boxed{\text{同支：}\ \frac1{|AF|}+\frac1{|BF|}=\frac{2a}{b^2}},\qquad
 \boxed{\text{异支：}\ \left|\frac1{|AF|}-\frac1{|BF|}\right|=\frac{2a}{b^2}}.
\]
同支看两点所在的分支，异支看焦点位于延长线上；所有距离均为正。

\end{statementbox}\clearpage\begin{statementbox}
  \textbf{\textcolor{CStatement}{核心结论}}
  \begin{description}[style=nextline,leftmargin=2.8em,labelsep=0.5em,itemsep=0.45em]
    \item[\textbf{结论一（同支：倒数相加）：}]
      当 $A,B$ 同支时，$F$ 位于线段 $AB$ 内，且
      \[
        \boxed{\frac1{|AF|}+\frac1{|BF|}=\frac{2a}{b^2}}.
      \]
    \item[\textbf{结论二（异支：倒数相减）：}]
      当 $A,B$ 异支时，$F$ 位于线段 $AB$ 的延长线上，且
      \[
        \boxed{\left|\frac1{|AF|}-\frac1{|BF|}\right|=\frac{2a}{b^2}}.
      \]
      \textbf{更强的有序式：} 若 $F=F_2(c,0)$，$A$ 在右支、$B$ 在左支，必有 $|AF|<|BF|$，故
      \[
        \boxed{\frac1{|AF|}-\frac1{|BF|}=\frac{2a}{b^2}}.
      \]
      左焦点时将“右支、左支”互换。
    \item[\textbf{统一常数（不用背两份）：}]
      设半通径参数 $p=\dfrac{b^2}{a}$、离心率 $e=\dfrac ca>1$，则
      \[
        \boxed{\frac{2a}{b^2}=\frac2p=\frac{2}{a(e^2-1)}}.
      \]
  \end{description}

  \medskip
  \textbf{\textcolor{CStatement}{范围提醒：}}
  斜率 $k=\pm\dfrac ba$ 的过焦点直线与一条渐近线平行，仅有一个有限交点，\textbf{不构成}上述两交点情形；竖直线 $x=\pm c$ 则有两个同支交点，\textbf{可以使用}结论。
\end{statementbox}
